% GENERATED FILE. Edit canonical lessons, then run npm run generate:books.
% canonical-source-hash: fnv1a64:733b6dff6e92a58a
% canonical-lessons: FA-C226-kud-dadan, FA-C226-paru-kardan, FA-C226-raha-kardan, FA-C226-sabidan, FA-C226-abkeshi-kardan

\chapter{To fertilise, To shovel, To let go, To scrub, To rinse}
\label{ch:fa-to-fertilise-to-shovel-to-let-go-to-scrub-to-rinse}
\hlchaptermodality{226}

\begin{chapteropening}
\textbf{By the end of this chapter:} \emph{I can say to fertilise, to shovel, to let go, to scrub and to rinse.}

Complete the last lesson of chapter 226: I can say to fertilise, to shovel, to let go, to scrub and to rinse.
\end{chapteropening}

\section[\texorpdfstring{\fa{کود} \fa{دادن} (kud dâdan)}{کود دادن (kud dâdan)}]{\fa{کود} \fa{دادن} (kud dâdan) — to fertilise}

\label{lesson:FA-C226-kud-dadan}



\begin{quote}
Before the new one: say the Persian for to vomit, then the Persian for to swell up.
\end{quote}

\subsection*{You'll want to know: \fa{کود} \fa{دادن}}
\textbf{\fa{کود} \fa{دادن}} — \emph{kud dâdan} — \textquotedblleft{}to fertilise\textquotedblright{}.

\begin{cousinweb}
\emph{kud dâdan} is \textquotedblleft{}to give manure\textquotedblright{}: \emph{kud}, \textquotedblleft{}manure, heap\textquotedblright{}, is a native word (EtymDB). \emph{dâdan}, \textquotedblleft{}to give\textquotedblright{}, is an inherited verb; EtymDB takes it back to the Indo-European root \emph{*deh\textsubscript{3}-}, \textquotedblleft{}to give\textquotedblright{}.
\end{cousinweb}

\begin{grammarlens}[title={What you've built}]
One more word for this chapter.
\end{grammarlens}

\subsection*{Your turn}
\begin{itemize}
\raggedright
  \item \emph{Say it:} \emph{kud dâdan}
  \item \emph{Say it:} \emph{kud dâdan}, once more
  \item \emph{Say it:} say \emph{kud dâdan}
  \item \emph{Recall:} say the Persian for to vomit, then the Persian for to swell up, then say \emph{kud dâdan} again
\end{itemize}

\begin{tcolorbox}[breakable,colback=teal!4,colframe=teal!35!black,title={Before you move on}]
What does \fa{کود} \fa{دادن} mean? (\textquotedblleft{}To fertilise\textquotedblright{}.) Say it once more.
\end{tcolorbox}
\section[\texorpdfstring{\fa{پارو} \fa{کردن} (pâru kardan)}{پارو کردن (pâru kardan)}]{\fa{پارو} \fa{کردن} (pâru kardan) — to shovel; to row}

\label{lesson:FA-C226-paru-kardan}



\begin{quote}
Before the new one: say the Persian for to swell up, then the Persian for to fertilise.
\end{quote}

\subsection*{You'll want to know: \fa{پارو} \fa{کردن}}
\textbf{\fa{پارو} \fa{کردن}} — \emph{pâru kardan} — \textquotedblleft{}to shovel; to row\textquotedblright{}.

\begin{cousinweb}
\emph{pâru kardan} pairs \emph{pâru}, \textquotedblleft{}shovel, oar\textquotedblright{}, a native Persian word, with \emph{kardan}, \textquotedblleft{}to do\textquotedblright{}. \emph{kardan}, \textquotedblleft{}to do\textquotedblright{}, is an inherited verb, from Middle Persian \emph{kardan} (EtymDB).
\end{cousinweb}

\begin{grammarlens}[title={What you've built}]
One more word for this chapter.
\end{grammarlens}

\subsection*{Your turn}
\begin{itemize}
\raggedright
  \item \emph{Say it:} \emph{pâru kardan}
  \item \emph{Say it:} \emph{pâru kardan}, once more
  \item \emph{Say it:} say \emph{pâru kardan}
  \item \emph{Recall:} say the Persian for to swell up, then the Persian for to fertilise, then say \emph{pâru kardan} again
\end{itemize}

\begin{tcolorbox}[breakable,colback=teal!4,colframe=teal!35!black,title={Before you move on}]
What does \fa{پارو} \fa{کردن} mean? (\textquotedblleft{}To shovel; to row\textquotedblright{}.) Say it once more.
\end{tcolorbox}
\section[\texorpdfstring{\fa{رها} \fa{کردن} (rahâ kardan)}{رها کردن (rahâ kardan)}]{\fa{رها} \fa{کردن} (rahâ kardan) — to let go, to release}

\label{lesson:FA-C226-raha-kardan}



\begin{quote}
Before the new one: say the Persian for to fertilise, then the Persian for to shovel.
\end{quote}

\subsection*{You'll want to know: \fa{رها} \fa{کردن}}
\textbf{\fa{رها} \fa{کردن}} — \emph{rahâ kardan} — \textquotedblleft{}to let go, to release\textquotedblright{}.

\begin{cousinweb}
\emph{rahâ kardan} pairs \emph{rahâ}, \textquotedblleft{}free, released\textquotedblright{}, with \emph{kardan}. \emph{rahâ} is native, usually traced to Middle Persian \emph{rahāg}. \emph{kardan}, \textquotedblleft{}to do\textquotedblright{}, is an inherited verb, from Middle Persian \emph{kardan} (EtymDB).
\end{cousinweb}

\begin{grammarlens}[title={What you've built}]
One more word for this chapter.
\end{grammarlens}

\subsection*{Your turn}
\begin{itemize}
\raggedright
  \item \emph{Say it:} \emph{rahâ kardan}
  \item \emph{Say it:} \emph{rahâ kardan}, once more
  \item \emph{Say it:} say \emph{rahâ kardan}
  \item \emph{Recall:} say the Persian for to fertilise, then the Persian for to shovel, then say \emph{rahâ kardan} again
\end{itemize}

\begin{tcolorbox}[breakable,colback=teal!4,colframe=teal!35!black,title={Before you move on}]
What does \fa{رها} \fa{کردن} mean? (\textquotedblleft{}To let go, to release\textquotedblright{}.) Say it once more.
\end{tcolorbox}
\section[\texorpdfstring{\fa{سابیدن} (sâbidan)}{سابیدن (sâbidan)}]{\fa{سابیدن} (sâbidan) — to scrub}

\label{lesson:FA-C226-sabidan}



\begin{quote}
Before the new one: say the Persian for to shovel, then the Persian for to let go.
\end{quote}

\subsection*{You'll want to know: \fa{سابیدن}}
\textbf{\fa{سابیدن}} — \emph{sâbidan} — \textquotedblleft{}to scrub\textquotedblright{}.

\begin{cousinweb}
\emph{sâbidan} is a native Persian verb, but no source used here gives its older form, and its origin is not settled.
\end{cousinweb}

\begin{grammarlens}[title={What you've built}]
One more word for this chapter.
\end{grammarlens}

\subsection*{Your turn}
\begin{itemize}
\raggedright
  \item \emph{Say it:} \emph{sâbidan}
  \item \emph{Say it:} \emph{sâbidan}, once more
  \item \emph{Say it:} say \emph{sâbidan}
  \item \emph{Recall:} say the Persian for to shovel, then the Persian for to let go, then say \emph{sâbidan} again
\end{itemize}

\begin{tcolorbox}[breakable,colback=teal!4,colframe=teal!35!black,title={Before you move on}]
What does \fa{سابیدن} mean? (\textquotedblleft{}To scrub\textquotedblright{}.) Say it once more.
\end{tcolorbox}
\section[\texorpdfstring{\fa{آبکشی} \fa{کردن} (âbkeshi kardan)}{آبکشی کردن (âbkeshi kardan)}]{\fa{آبکشی} \fa{کردن} (âbkeshi kardan) — to rinse}

\label{lesson:FA-C226-abkeshi-kardan}



\begin{quote}
Before the new one: say the Persian for to let go, then the Persian for to scrub.
\end{quote}

\subsection*{You'll want to know: \fa{آبکشی} \fa{کردن}}
\textbf{\fa{آبکشی} \fa{کردن}} — \emph{âbkeshi kardan} — \textquotedblleft{}to rinse\textquotedblright{}.

\begin{cousinweb}
\emph{âbkeshi kardan} pairs \emph{âbkeshi}, \textquotedblleft{}rinsing\textquotedblright{}, with \emph{kardan}. \emph{âbkeshi} is \emph{âb}, \textquotedblleft{}water\textquotedblright{} (EtymDB), with \emph{kesh}, \textquotedblleft{}drawing\textquotedblright{}, and \emph{-i}. All are native Persian. \emph{kardan}, \textquotedblleft{}to do\textquotedblright{}, is an inherited verb, from Middle Persian \emph{kardan} (EtymDB).
\end{cousinweb}

\begin{grammarlens}[title={What you've built}]
One more word for this chapter.
\end{grammarlens}

\subsection*{Your turn}
\begin{itemize}
\raggedright
  \item \emph{Say it:} \emph{âbkeshi kardan}
  \item \emph{Say it:} \emph{âbkeshi kardan}, once more
  \item \emph{Say it:} say \emph{âbkeshi kardan}
  \item \emph{Recall:} say the Persian for to let go, then the Persian for to scrub, then say \emph{âbkeshi kardan} again
\end{itemize}

\begin{tcolorbox}[breakable,colback=teal!4,colframe=teal!35!black,title={Before you move on}]
What does \fa{آبکشی} \fa{کردن} mean? (\textquotedblleft{}To rinse\textquotedblright{}.) Say it once more.
\end{tcolorbox}
